% f2015gdw_q17_histogram.tex — 2015 年高考数学·广东卷·文科 第 17 题（频率分布直方图）
%
% 【来源】`.work/phase3-b16/source/2015广东文.pdf`（SHA-256
%   0c27f1a08f837014066c636a3190b8fa7ad7e5a4494f204e60412293aea17e09），
%   第 1 页（合订本页 949）第 3 栏；内容对象 5。逐路径实测见
%   `.work/phase3-b16/figure-inventory.md` Q17 节。
%
% 【坐标口径】清单坐标为「真页 pt、原点左上、y 向下」。本图取
%   X0 = 862.9140（y 轴 x），Y0 = 345.5480（x 轴基线 y_top），
%   本地坐标  (x, y) = (x_page − X0, Y0 − y_top)，
%   即本地原点 = 坐标轴交点，x 向右、y 向上，1 单位 = 1 TeX pt。
%   换算示例：原点 (862.9140,345.5479) → (0, 0.0001)；
%   柱 [160,180) 左边界 (885.5915,345.5485) → (22.6775,−0.0005)；
%   y 轴箭头尖 (862.9140,127.1584) → (0,218.3896)。
%   y 轴箭头的旋转坐标系按 PDF 行向量右乘（后写 cm 先作用）换算（FI-S5）。
%
% 【线宽/虚线】轴线与柱体 line width = 0.398pt；分式横线 0.42pt；
%   两枚箭头为填充路径（源 `f`，无描边；源 lw=1.0，填充路径下线宽不参与着色），
%   包围盒 2.99 × 3.99pt；依赖 \usetikzlibrary{arrows.meta}（按库规加载；
%   箭头用路径级贝塞尔复刻，未用 Triangle 造型）；
%   虚线导引线 dash pattern = on 2.989pt off 2.989pt，dash phase = 0pt。
%
% 【照源复刻的存疑项】
%   FI-S4：x 导引线 y_top=228.6168，[240,260) 柱顶 y_top=228.5026，
%          相差 0.1142pt；照录（导引线与柱顶分别定位）。
%   FI-S5：y 轴箭头在 `0 1 -1 0 862.914 711.14509 cm` 旋转坐标系内，
%          按行向量右乘解析，箭头包围盒 x 861.42–864.41、y_top 127.16–131.14。
%   FI-S8：底部 0.0025 / 0.002 基线仅差 7.85pt，照录。
%   FI-S10：x 轴 160…300 八个刻度在源内为单个 TJ 以字距排开，
%          此处按逐点 Tm 起点分别定位（起点取自 pdftotext -bbox）。
%   断轴三齿折线逐点照录（11 点；振幅 ±2.339pt）。
%
% 【自验结果】见文件末尾。
%
% 元素计数：18 条路径（轴+断轴 1、箭头 2、分式线 1、柱 7、虚线导引 7）、
%   21 个文本标签。
\begin{tikzpicture}[x=1pt, y=1pt,
    every node/.style={inner sep=0pt, font=\fontsize{10.5pt}{12.6pt}\selectfont}]

  % ---- y 轴 + 断轴 + x 轴（源 path 26，14 点折线，lw 0.398） ----
  % 源真页坐标：m (862.9140,130.7449) l (862.9140,345.5479) l (867.4490,345.5479)
  %   l (868.5830,343.2099) l (870.8510,347.8869) l (871.9850,345.5479)
  %   l (873.1190,343.2099) l (875.3870,347.8869) l (876.5210,345.5489)
  %   l (877.6550,343.2099) l (879.9230,347.8869) l (881.0570,345.5489)
  %   l (885.5920,345.5489) l (1086.1030,345.5479)
  \draw[line width=0.398pt]
    (0.0000,214.8031) -- (0.0000,0.0001)
    -- (4.5350,0.0001) -- (5.6690,2.3381) -- (7.9370,-2.3389)
    -- (9.0710,0.0001) -- (10.2050,2.3381) -- (12.4730,-2.3389)
    -- (13.6070,-0.0009) -- (14.7410,2.3381) -- (17.0090,-2.3389)
    -- (18.1430,-0.0009) -- (22.6780,-0.0009) -- (223.1890,0.0001);

  % ---- y 轴箭头（源 path 27，两段贝塞尔，填充，无描边） ----
  % 源真页坐标（FI-S5 校正后）：尖 (862.9140,127.1584)
  %   c (862.7150,128.2214) (862.1170,129.9484) (861.4200,131.1434)
  %   l (864.4080,131.1434)
  %   c (863.7110,129.9484) (863.1130,128.2214) (862.9140,127.1584)
  \fill[line width=1.0pt]
    (0.0000,218.3896)
    .. controls (-0.1990,217.3266) and (-0.7970,215.5996) .. (-1.4940,214.4046)
    -- (1.4940,214.4046)
    .. controls (0.7970,215.5996) and (0.1990,217.3266) .. (0.0000,218.3896)
    -- cycle;

  % ---- x 轴箭头（源 path 28，两段贝塞尔，填充，无描边） ----
  % 源真页坐标：尖 (1089.6893,345.5480)
  %   c (1088.6263,345.3490) (1086.8993,344.7510) (1085.7043,344.0540)
  %   l (1085.7043,347.0420)
  %   c (1086.8993,346.3450) (1088.6263,345.7470) (1089.6893,345.5480)
  \fill[line width=1.0pt]
    (226.7753,0.0000)
    .. controls (225.7123,0.1990) and (223.9853,0.7970) .. (222.7903,1.4940)
    -- (222.7903,-1.4940)
    .. controls (223.9853,-0.7970) and (225.7123,-0.1990) .. (226.7753,0.0000)
    -- cycle;

  % ---- 「频率/组距」分式横线（源 path 29，lw 0.42） ----
  % 源真页坐标：(867.8080,127.8390) → (888.8080,127.8390)
  \draw[line width=0.42pt] (4.8940,217.7090) -- (25.8940,217.7090);

  % ---- 柱 [160,180)：源 path 30，m,l,l,l,h，lw 0.398 ----
  \draw[line width=0.398pt]
    (22.6775,31.1645) -- (22.6775,-0.0005) -- (45.3555,-0.0005)
    -- (45.3555,31.1645) -- cycle;
  % ---- 导引线 0.002：源 path 31，虚线 ----
  \draw[line width=0.398pt, dash pattern=on 2.989pt off 2.989pt, dash phase=0pt]
    (0.0000,31.1645) -- (22.6780,31.1645);

  % ---- 柱 [180,200)：源 path 32 ----
  \draw[line width=0.398pt]
    (45.3551,148.2099) -- (45.3551,-0.0001) -- (68.0331,-0.0001)
    -- (68.0331,148.2099) -- cycle;
  % ---- 导引线 0.0095：源 path 33 ----
  \draw[line width=0.398pt, dash pattern=on 2.989pt off 2.989pt, dash phase=0pt]
    (0.0000,148.2099) -- (45.3550,148.2099);

  % ---- 柱 [200,220)：源 path 34 ----
  \draw[line width=0.398pt]
    (68.0326,171.5239) -- (68.0326,-0.0001) -- (90.7106,-0.0001)
    -- (90.7106,171.5239) -- cycle;
  % ---- 导引线 0.011：源 path 35 ----
  \draw[line width=0.398pt, dash pattern=on 2.989pt off 2.989pt, dash phase=0pt]
    (0.0000,171.5239) -- (68.0330,171.5239);

  % ---- 柱 [220,240)：源 path 36 ----
  \draw[line width=0.398pt]
    (90.7101,194.8378) -- (90.7101,-0.0002) -- (113.3881,-0.0002)
    -- (113.3881,194.8378) -- cycle;
  % ---- 导引线 0.0125：源 path 37 ----
  \draw[line width=0.398pt, dash pattern=on 2.989pt off 2.989pt, dash phase=0pt]
    (0.0000,194.8378) -- (90.7100,194.8378);

  % ---- 柱 [240,260)：源 path 38（注意本柱顶无独立导引线，见 FI-S4） ----
  \draw[line width=0.398pt]
    (113.3877,117.0454) -- (113.3877,0.0004) -- (136.0657,0.0004)
    -- (136.0657,117.0454) -- cycle;

  % ---- 柱 [260,280)：源 path 39 ----
  \draw[line width=0.398pt]
    (136.0652,78.0303) -- (136.0652,0.0003) -- (158.7432,0.0003)
    -- (158.7432,78.0303) -- cycle;
  % ---- 导引线 0.005：源 path 40 ----
  \draw[line width=0.398pt, dash pattern=on 2.989pt off 2.989pt, dash phase=0pt]
    (0.0000,78.0303) -- (136.0650,78.0303);

  % ---- 柱 [280,300)：源 path 41 ----
  \draw[line width=0.398pt]
    (158.7427,39.0151) -- (158.7427,0.0001) -- (181.4207,0.0001)
    -- (181.4207,39.0151) -- cycle;
  % ---- 导引线 0.0025：源 path 42 ----
  \draw[line width=0.398pt, dash pattern=on 2.989pt off 2.989pt, dash phase=0pt]
    (0.0000,39.0151) -- (158.7430,39.0151);

  % ---- 导引线 x：源 path 43（y_top 228.6168，与柱顶差 0.1142pt，FI-S4） ----
  \draw[line width=0.398pt, dash pattern=on 2.989pt off 2.989pt, dash phase=0pt]
    (0.0000,116.9312) -- (113.3880,116.9312);

  % ---- 文本标签：源 Tm/Td 基线原点，10.5pt ----
  % 「频率/组距」：STSong；源基线 (867.8080,123.3610) / (867.8080,137.6670)
  \node[anchor=base west] at (4.8940,222.1870) {频率};
  \node[anchor=base west] at (4.8940,207.8810) {组距};
  % y 轴刻度值（源 C0_1 LMRoman10；基线左端）
  \node[anchor=base west] at (-32.8680,191.4570) {$0.0125$};  % 源 (830.0460,154.0910)
  \node[anchor=base west] at (-27.6180,168.1430) {$0.011$};   % 源 (835.2960,177.4050)
  \node[anchor=base west] at (-32.8680,144.8290) {$0.0095$};  % 源 (830.0460,200.7190)
  \node[anchor=base west] at (-9.7000,114.6710)  {$x$};       % 源 (853.2140,230.8770) CMMI10
  \node[anchor=base west] at (-27.6180,74.6490)  {$0.005$};   % 源 (835.2960,270.8990)
  \node[anchor=base west] at (-32.8680,35.6340)  {$0.0025$};  % 源 (830.0460,309.9140)
  \node[anchor=base west] at (-27.6180,27.7830)  {$0.002$};   % 源 (835.2960,317.7650)
  % 原点 0（源 (860.2890,356.2400)）
  \node[anchor=base west] at (-2.6250,-10.6920) {$0$};
  % x 轴刻度（源为单个 TJ；以下为逐点 Tm 起点，基线 y_top=356.0140）
  \node[anchor=base west] at (14.8020,-10.4660)  {$160$};  % 源 (877.716000,356.014)
  \node[anchor=base west] at (37.4800,-10.4660)  {$180$};  % 源 (900.394005,356.014)
  \node[anchor=base west] at (60.1570,-10.4660)  {$200$};  % 源 (923.071002,356.014)
  \node[anchor=base west] at (82.8350,-10.4660)  {$220$};  % 源 (945.749007,356.014)
  \node[anchor=base west] at (105.5130,-10.4660) {$240$};  % 源 (968.427012,356.014)
  \node[anchor=base west] at (128.1900,-10.4660) {$260$};  % 源 (991.104009,356.014)
  \node[anchor=base west] at (150.8680,-10.4660) {$280$};  % 源 (1013.782014,356.014)
  \node[anchor=base west] at (173.5450,-10.4660) {$300$};  % 源 (1036.459011,356.014)
  % x 轴标题「月平均用电量/度」：源三段分别定位，同一基线 (…,339.2240)
  \node[anchor=base west] at (187.4000,6.3240) {月平均用电量};  % 源 (1050.3140,339.2240) STSong
  \node[anchor=base west] at (250.4000,6.3240) {/};             % 源 (1113.3140,339.2240) LMRoman10
  \node[anchor=base west] at (255.6500,6.3240) {度};            % 源 (1118.5640,339.2240) STSong

\end{tikzpicture}
%
% 【自验结果】
%   探针：latexmk -pdfxe，0 error / 0 Overfull / 0 Underfull / 0 Missing character；
%         页面 301.40 × 240.11 pt（standalone border=0）。
%   tex 反算：18 条路径 + 21 个标签全部匹配，最大残差 0.000020 pt（阈值 0.05）。
%   600dpi 墨迹：整体墨迹包围盒最大边差 0.278 pt；几何包围盒 0.213 pt；
%         标签中心最大 0.273 pt（280）。（尺度归一，已扣除 0.996277 的
%         TeX pt / PDF bp 均匀因子）
%         按页面 bp 直接比较（物理）最大 0.952 pt（几何包围盒），来自该 0.37% 均匀缩放。
%   灵敏度：y 轴线宽 +0.02pt → 600dpi 等效线宽 +0.0358（tex 检查器不检线宽，报 0）；
%         刻度 200 y +0.15pt → tex 检查器报 0.1500、600dpi 标签中心报 0.2409。均能报出。
